% !TEX root = ../algebra_02.tex

\section{Матрица линейного отображения. Матрица композиции}

\begin{Definition}{Матрица линейного отображения}{}
	Пусть $\alpha\colon V\to W$ --- линейное отображение. Его матрицей в базисах $E$ (пространства $V$) и $F$ (пространства $W$) называется матрица $[\alpha]_{E,F}\in M_{m\times n}(K)$, столбцы которой равны координатам образов базисных векторов:
	\[
	[\alpha]_{E,F} = \bigl([\alpha e_1]_F,\ldots,[\alpha e_n]_F\bigr)
	\].
	Иначе говоря, выполняется символическое равенство $\alpha E=F[\alpha]_{E,F}$.
\end{Definition}

Матрица композиции линейных отображений тесно связана с матричным умножением.

\begin{Theorem}{Матрица композиции}{}
	Пусть заданы линейные отображения
	\[
	V \xrightarrow{\alpha} W \xrightarrow{\beta} U,
	\]
	а также базисы $E$ в $V$, $F$ в $W$ и $G$ в $U$.
	Для композиции $\beta\alpha\colon V\to U$ матрица в соответствующих базисах равна произведению матриц этих отображений:
	\[
	[\beta\alpha]_{E,G} = [\beta]_{F,G}[\alpha]_{E,F}
	\].
\end{Theorem}

\begin{proof}
	Пусть $[\alpha]_{E,F}=A$ и $[\beta]_{F,G}=B$. По определению матрицы линейного отображения имеем $\alpha E=FA$ и $\beta F=GB$.
	Применим композицию к базису $E$:
	\[
	(\beta\alpha)E = \beta(\alpha E) = \beta(FA) = (\beta F)A = (GB)A = G(BA)
	\].
	Следовательно, матрица композиции $[\beta\alpha]_{E,G}$ равна $BA$, то есть $[\beta]_{F,G}[\alpha]_{E,F}$.
\end{proof}
